\documentclass[11pt,a4paper]{article}
\usepackage[margin=1in]{geometry}
\usepackage[T2A]{fontenc}
\usepackage[utf8]{inputenc}
\usepackage[english,russian]{babel}
\usepackage{amsmath,amssymb}
\usepackage{microtype}
\usepackage{hyperref}

\setlength{\parindent}{0pt}
\setlength{\parskip}{0.75\baselineskip}
\sloppy

\title{Горизонт как поверхность фазовой записи в модели 01\\\large Чёрная дыра без точечной сингулярности и без буквального испарения}
\author{}
\date{}

\newcommand{\field}{\mathcal{F}_{01}}
\newcommand{\phase}{\varphi}
\newcommand{\normal}{\mathcal{N}}
\newcommand{\memory}{\mathcal{M}}
\newcommand{\bulk}{\mathrm{bulk}}
\newcommand{\boundary}{\mathrm{boundary}}

\begin{document}
\maketitle

\begin{center}
\textit{Рабочая версия v0.2 --- русская исследовательская редакция}
\end{center}

\begin{abstract}
В статье формулируется отдельный горизонтный блок модели 01. Чёрная дыра рассматривается не как объект с точечной физической сингулярностью внутри, а как предельный режим поля, в котором объёмная фазовая структура переходит в поверхностную запись. В этом языке горизонт является не входом к бесконечной плотности, а границей, где исчезает радиальная нормаль, прекращается объёмное разворачивание динамики и сохраняется фазовая память. Энтропия Бекенштейна--Хокинга интерпретируется как указание на поверхностный характер записи. При этом модель 01 принципиально отличается от буквальной картины испарения чёрной дыры: горизонт не теряет запись как тепловой излучатель, а накапливает фазовую информацию. Такой подход пока является гипотезой и требует строгого сопоставления с полуклассической теорией Хокинга, квантовой теорией поля в искривлённом пространстве-времени и голографическими подходами.
\end{abstract}

\textbf{Ключевые слова:} модель 01; чёрная дыра; горизонт; фазовая запись; память; Бекенштейн--Хокинг; голография; сингулярность; излучение Хокинга; квантовая теория поля.

\section{Как читать эту статью}

Эта статья является продолжением работы об элементарной частице как замкнутой волне поля 01. В первой статье локальная частица описывалась как замкнутый фазовый узел с нормалью, массой и памятью. Здесь рассматривается предельный случай: что происходит, когда такая локальная структура достигает горизонта чёрной дыры.

Текст следует читать на трёх уровнях:
\begin{enumerate}
\item \textbf{Известная физика}: в общей относительности чёрные дыры обладают горизонтом; в термодинамике чёрных дыр энтропия пропорциональна площади горизонта; в полуклассической теории Хокинга горизонт связан с термальным излучением.
\item \textbf{Интерпретация модели 01}: горизонт понимается как поверхность записи, где объёмная фазовая структура переводится в поверхностный код.
\item \textbf{Открытая гипотеза}: модель 01 не принимает буквальное испарение чёрной дыры как потерю записи через тепловой поток; вместо этого предполагается накопление фазовой информации на горизонте.
\end{enumerate}

Поэтому статья не утверждает, что стандартные расчёты Хокинга уже опровергнуты. Она формулирует альтернативный язык, который должен быть отдельно сопоставлен с полуклассической физикой горизонтов.

\section{Что принимается и что переинтерпретируется}

Чтобы избежать неверного чтения, зафиксируем границу между стандартной физикой и моделью 01. Модель не отрицает саму роль горизонта, энтропии и редуцированного описания. Она переинтерпретирует физический смысл этих структур.

\begin{center}
\begin{tabular}{p{0.43\textwidth}p{0.43\textwidth}}
\textbf{Стандартное описание} & \textbf{Интерпретация модели 01} \\
горизонт ограничивает причинный доступ & горизонт является поверхностью фазовой записи \\
энтропия пропорциональна площади & площадь измеряет ёмкость поверхностного архива \\
внешний наблюдатель видит редуцированное состояние & внешнему наблюдателю доступна только часть фазовой записи \\
излучение Хокинга имеет термальный характер & термальность читается как эффект ограниченного доступа \\
испарение описывает потерю массы через поток & буквальная потеря записи не принимается; запись накапливается на горизонте \\
сингулярность возникает в классическом продолжении внутрь & точечный центр заменяется предельной поверхностной фиксацией.
\end{tabular}
\end{center}

Эта таблица не является доказательством модели. Она задаёт карту различий: что модель 01 принимает как важный физический факт, а что переводит на свой язык.


\section{Главная идея}

В модели 01 чёрная дыра не понимается как дырка с физической точкой бесконечной плотности внутри. Она понимается как предельный режим поля:
\begin{equation}
\text{ЧД} \neq \text{точечная сингулярность внутри},
\qquad
\text{ЧД} \sim \text{поверхность предельной фиксации}.
\end{equation}

Ключевая формулировка модели:
\begin{equation}
\text{ноль не находится внутри; ноль и есть граница.}
\end{equation}

Это означает, что предельный режим не помещается в центр как маленькая область бесконечной плотности. Напротив, само понятие внутреннего объёма перестаёт быть главным. Физически значимым объектом становится горизонт как поверхность, на которой фиксируется фазовая структура.

\section{Локальная память частицы и глобальная запись горизонта}

В модели 01 элементарная частица является локальным узлом памяти. Она удерживает фазовую структуру через замкнутый обход и нормаль. Схематически:
\begin{equation}
\text{ЭЧ} = \text{локальная память} + \normal.
\end{equation}

Горизонт чёрной дыры является другим режимом памяти. Он не удерживает локальную нормаль отдельной частицы, а переводит множество локальных фазовых структур в глобальную поверхностную запись:
\begin{equation}
\memory_{\mathrm{particle}} \longrightarrow \memory_{\mathrm{horizon}}.
\end{equation}

В терминах объёма и границы этот переход можно записать так:
\begin{equation}
\memory_{\bulk} \longrightarrow \memory_{\boundary}.
\end{equation}

Эта запись пока не является строгим оператором отображения. Она фиксирует принцип, который должен быть формализован: фазовая информация, связанная с объёмной структурой, должна быть представлена на поверхности горизонта.

\section{Связь с энтропией Бекенштейна--Хокинга}

Формула энтропии чёрной дыры имеет вид
\begin{equation}
S_{\mathrm{BH}} = \frac{k_B A}{4\ell_P^2},
\end{equation}
где $A$ --- площадь горизонта, $\ell_P$ --- планковская длина, $k_B$ --- постоянная Больцмана.

Для стандартной физики эта формула показывает глубокую связь между горизонтом, термодинамикой, квантовой теорией и гравитацией. Для модели 01 важен следующий факт: энтропия масштабируется с площадью, а не с объёмом. В языке модели это читается как указание на поверхностный характер записи:
\begin{equation}
\text{число доступных записей} \sim \frac{A}{\ell_P^2}.
\end{equation}

Модель 01 не выводит эту формулу в строгом виде. Она предлагает интерпретацию: площадь горизонта измеряет ёмкость поверхностного фазового архива.

\section{Почему точечная сингулярность не требуется}

Классическая проблема сингулярности возникает, если считать, что всё падающее вещество должно продолжать движение внутрь и в итоге сжаться в точку. Модель 01 предлагает иной сценарий.

При приближении к горизонту локальная структура теряет радиальную нормаль:
\begin{equation}
\normal \to 0.
\end{equation}

Если нормаль исчезает, исчезает и локальная объёмная глубина. Тогда структура не обязана продолжать разворачиваться внутрь. Она может перейти в поверхностный режим:
\begin{equation}
\text{объёмная динамика} \longrightarrow \text{поверхностная запись}.
\end{equation}

В этом смысле модель заменяет схему
\begin{equation}
\text{вещество} \longrightarrow \text{точка бесконечной плотности}
\end{equation}
на схему
\begin{equation}
\text{фазовая структура} \longrightarrow \text{горизонтная запись}.
\end{equation}

Поэтому сингулярность не устраняется утверждением, что внутри находится новая неизвестная материя. Она устраняется сменой описания: центр перестаёт быть местом хранения информации, а главным становится горизонт.

\section{Что происходит с замкнутой волной у горизонта}

Локальная частица в модели 01 имеет тангенциальные обходы и нормаль. При приближении к предельному режиму можно выделить четыре шага:
\begin{enumerate}
\item тангенциальные фазовые обходы вытягиваются вдоль горизонта;
\item радиальная нормаль подавляется;
\item объёмная глубина узла исчезает;
\item фазовый код переписывается в поверхностную структуру.
\end{enumerate}

Схематически:
\begin{equation}
\left(1_{\tau_1},1_{\tau_2}\right)+\normal_0
\quad \longrightarrow \quad
\text{поверхностный фазовый код}.
\end{equation}

Это означает, что частица не падает в физическую точку. Она теряет нормаль и перестаёт быть локальным объёмным узлом. Её память переходит в горизонтный режим.

\section{Отличие от буквального испарения}

В стандартной полуклассической картине излучение Хокинга описывает термальный поток, связанный с квантовыми полями на фоне горизонта. Эта картина является важной частью современной теоретической физики и не должна подменяться простым отрицанием.

В версии v0.2 позиция модели 01 формулируется осторожно: расчёт Хокинга рассматривается как описание термального отклика внешнего наблюдателя на горизонт и редуцированные степени свободы. Но модель не принимает дополнительную буквальную картину, в которой сама чёрная дыра является обычным телом, постепенно теряющим внутреннюю запись как тепло. В её языке чёрная дыра не является объектом, который буквально испаряется и уничтожает фазовую память через тепловой поток:
\begin{equation}
\begin{aligned}
\text{ЧД} &\neq \text{испаряющийся термодинамический объект},\\
\text{ЧД} &\sim \text{растущая поверхность фазовой записи}.
\end{aligned}
\end{equation}

То, что в стандартном описании выглядит как термальность, в модели 01 должно быть переосмыслено как следствие ограниченного доступа к фазовой информации. Наблюдатель вне горизонта имеет доступ не к полной фазовой структуре, а к редуцированному описанию:
\begin{equation}
\rho_{\mathrm{out}} = \operatorname{Tr}_{\mathrm{in}}\left(|\Psi\rangle\langle\Psi|\right).
\end{equation}

В стандартном языке такая редукция связана с энтропией и термальными свойствами. В языке модели 01 она указывает на неполноту внешнего доступа к записи, а не обязательно на потерю самой записи.

Важно подчеркнуть: это сильное отличие модели 01. Поэтому его нельзя подавать как доказанный факт. Корректная формулировка такова: модель 01 выдвигает гипотезу, что горизонт является архивом фазовой информации, а не механизмом окончательного теплового исчезновения записи. Следовательно, будущая работа должна показать, можно ли получить термальный спектр как эффект редуцированного доступа, не вводя буквальную потерю фазового архива.

\section{Горизонт как фазовый архив}

Если горизонт является поверхностью записи, то его рост означает не только увеличение массы или площади в обычном смысле. Он означает увеличение ёмкости архива:
\begin{equation}
A \uparrow \quad \Rightarrow \quad \dim \mathcal{H}_{\mathrm{horizon}} \uparrow.
\end{equation}

Здесь $\mathcal{H}_{\mathrm{horizon}}$ обозначает не строго построенное гильбертово пространство, а будущий кандидат на пространство состояний горизонта. Формула показывает направление формализации: рост площади должен соответствовать росту числа различимых фазовых записей.

В таком языке информация не уничтожается в центре и не уносится наружу как обычное тепло. Она меняет представление:
\begin{equation}
\text{локальный узел} \longrightarrow \text{поверхностный код}.
\end{equation}

\section{Последняя замкнутая волна и цикл}

В книге вводится предельный образ последней элементарной частицы или последней замкнутой волны. Это не утверждение о конкретной наблюдаемой частице, а космологический образ предельного различия.

Пока существует хотя бы один замкнутый узел, сохраняются:
\begin{itemize}
\item различие между режимами;
\item направление $1$;
\item минимальная фазовая задержка;
\item возможность счёта изменений, то есть время.
\end{itemize}

Когда последняя замкнутая волна переходит в режим записи, исчезает не объект в обычном смысле, а последняя поддерживаемая фазовая задержка:
\begin{equation}
\text{последняя ЭЧ} \longrightarrow \text{последняя горизонтная запись}.
\end{equation}

После этого внешний режим больше не удерживает различия. В книге это описывается как предельный переход $0$--$0$: совпадение двух режимов фиксации, из которого может возникнуть новый минимальный импульс различия.

Этот раздел пока остаётся наиболее спекулятивным. Его задача --- сохранить связь с космологической частью книги, но не смешивать её с более осторожной физической частью статьи о горизонте.

\section{Что нужно формализовать}

Чтобы горизонтная часть модели 01 стала физической теорией, необходимо определить:
\begin{enumerate}
\item пространство состояний горизонта $\mathcal{H}_{\mathrm{horizon}}$;
\item отображение $\memory_{\bulk}\to\memory_{\boundary}$;
\item закон сохранения фазовой информации;
\item связь площади горизонта с числом различимых фазовых записей;
\item отношение модели к расчёту Хокинга;
\item наблюдательные или теоретические отличия от стандартной картины испарения;
\item связь с голографическим принципом и black hole information problem.
\end{enumerate}

До выполнения этих шагов модель 01 должна оставаться осторожной: она предлагает язык и гипотезу, а не завершённое доказательство.

\section{Заключение}

Модель 01 предлагает рассматривать чёрную дыру как предельный режим поля, где локальные фазовые структуры теряют нормаль и переходят в поверхностную запись. В этом описании горизонт становится не входом к точечной сингулярности, а архивом фазовой информации.

Главное отличие от буквального испарения состоит в следующем: горизонт не уничтожает и не теряет запись как обычный тепловой излучатель. Он ограничивает внешний доступ к информации и переводит объёмную структуру в поверхностный код. Это согласуется с мотивом площади в формуле Бекенштейна--Хокинга, но требует отдельной строгой формализации и сопоставления с полуклассическими расчётами.

Следующий шаг --- построить математическое отображение между объёмной памятью и поверхностной записью, а также уточнить, каким образом термальность горизонта возникает как редуцированное описание, а не как буквальное исчезновение фазового архива.

\end{document}